\section{Introduction \& Research Question}
Urban green infrastructure, such as tree canopy cover and public parks, is widely recognized for its environmental and psychological benefits. However, its economic valuation is often implicit rather than direct. In urban economics, the hedonic pricing method (Rosen, 1974) posits that the price of a heterogeneous good, such as residential property, is a function of its structural characteristics (e.g., size, age) and neighborhood environmental amenities. 

While public investment in urban forestry is frequently justified by carbon sequestration and stormwater management, quantifying the capitalization of green space into housing prices provides crucial empirical evidence for municipal budgeting and local tax planning. 

This study addresses the following primary research question:
\begin{quote}
\textit{To what extent does neighborhood green space cover influence residential property values, holding physical property characteristics constant?}
\end{quote}

We test the null hypothesis that neighborhood green space has no effect on property value ($H_0: \beta_1 = 0$) against the alternative hypothesis that it has a positive effect ($H_1: \beta_1 > 0$).
